\section{Related Work}
\label{sec:related-work}

Regression-based ocular correction estimates how eye potentials propagate to each scalp channel and subtracts the predicted contribution \cite{gratton1983ocular,semlitsch1986ocular}. Its correction coefficients have a direct physical reading, and later work showed that eye-related subspaces estimated once from a calibration recording can be applied on-line with matrix operations \cite{kobler2020sgeyesub}. Blind source separation and artifact subspace reconstruction instead exploit the spatial structure of multichannel EEG and need no EOG channel \cite{jung2000bss,chang2020asr}. This literature established both the value of a calibration segment and the linear propagation model we build on, but it treats clean EEG only implicitly, as whatever remains after subtraction. Our method keeps the calibrated propagation model of this lineage and adds what it lacks, a learned distribution of clean EEG that decides how the subtraction should be corrected.

Supervised networks trained on semi-simulated corpora such as EEGdenoiseNet learn flexible mappings from contaminated to clean signals \cite{zhang2021eegdenoisenet,sun2020rescnn,yu2022deepseparator,chuang2022icunet,lai2024cleegn}, and calibration-free models extend this to unseen montages \cite{marcosmartinez2025eegoarnet}. In these models participant variation is absorbed into shared parameters. The closest prior work to ours in intent is the domain-specific denoising diffusion model of \citet{duan2023dsddpm}, which makes diffusion-based EEG denoising subject-aware by learning subject-specific noise distributions inside the network, so that the subjects to be denoised participate in training. We share the goal and the diffusion backbone but place the subject information elsewhere: the network stays subject-independent, and subject specificity enters through a closed-form propagation estimate from the evaluated user's own calibration segment. The two designs answer different deployment questions. Learning subject structure into the weights suits a fixed user pool; measuring it at calibration time extends to users the network has never seen, which is the setting we evaluate.

Diffusion models learn to invert a prescribed noising process \cite{ho2020ddpm,song2021scoresde} and, conditioned on a degraded observation, restore images and audio \cite{saharia2022palette,richter2023storm}; EEGDfus and D4PM apply conditional diffusion to EEG artifact removal \cite{huang2025eegdfus,shao2025d4pm}. When the degradation operator is known, restoration methods such as DDRM and diffusion posterior sampling use it explicitly during sampling \cite{kawar2022ddrm,chung2023dps}; when it is unknown, blind approaches estimate operator and signal jointly \cite{chung2023blinddps,murata2023gibbsddrm}. Our setting sits between these cases. The operator is unknown a priori but measurable: a calibration segment provides it in closed form before sampling begins, with an uncertainty that we propagate into the sample distribution. We evaluate the resulting intervals by empirical coverage and the continuous ranked probability score \cite{gneiting2007strictly}, rather than treating sampler variability as calibrated uncertainty by construction.
